% LaTeX Canvas: Chapter Outline
\chapter{Text Editors}
\label{ch:text-editors}

\section*{Purpose}
Text editors are the primary tool for writing code, configuring tools, inspecting logs, and fixing small problems quickly. For beginners, the main challenge is not typing---it is developing reliable workflows: opening the right file, editing safely, finding and replacing correctly, and understanding the trade-offs between terminal editors, GUI editors, and full IDEs.

\section*{Learning objectives}
By the end of this chapter, a reader should be able to:
\begin{enumerate}
\item Explain what makes a file ``text'' (vs binary) and why encoding and line endings matter.
\item Use essential editor operations: open/save, undo/redo, find/replace, multi-file search, and navigation.
\item Debug common file problems: wrong path, wrong extension, hidden characters, encoding issues, and line-ending mismatches.
\item Choose an appropriate editor class for a task (terminal editor vs GUI editor vs IDE).
\item Use a GUI editor for simple scripting: run scripts, read errors, and iterate.
\item Apply safe refactoring habits with find/replace and version control.
\item Configure an editor minimally: indentation, formatting, linting, and extensions.
\end{enumerate}

\section*{Running theme: editors are tools, not identities}
Choose an editor that fits the task, and build habits that transfer across editors: paths, file formats, search, safe edits, and reproducibility.

\section{A beginner mental model}
\subsection{Text files are infrastructure}
\begin{itemize}
\item Code, configuration, logs, and data dictionaries are often plain text.
\item Editors are the "workbench" for these artifacts.
\end{itemize}

\subsection{What an editor does (conceptually)}
\begin{itemize}
\item Reads bytes from disk and interprets them as text using an \textbf{encoding} (often UTF-8).
\item Displays text with line breaks and whitespace.
\item Writes bytes back to disk when you save.
\end{itemize}

\subsection{Key terms}
\begin{description}
\item[Encoding] How bytes become characters (UTF-8, etc.).
\item[Line endings] Newline conventions (LF vs CRLF) that affect cross-platform projects.
\item[Whitespace] Spaces/tabs/newlines that can change program behavior.
\item[Syntax highlighting] Visual cues based on language rules.
\item[Linting/formatting] Automated checks and style normalization.
\item[IDE] Integrated Development Environment: editor + build + run + debug + tooling.
\end{description}

\section{Choosing your tool: three editor classes}
\subsection{Terminal editors (nano, vim, emacs)}
\begin{itemize}
\item Best for: remote servers, quick config edits, minimal environments.
\item Strengths: always available, low overhead, works over SSH.
\item Trade-offs: steeper learning curve (especially modal editing), fewer visual affordances.
\item Recommended student baseline: learn \textbf{nano} for emergencies; optionally learn \textbf{vim} motions over time.
\end{itemize}

\subsection{GUI code editors (VS Code, Sublime, etc.)}
\begin{itemize}
\item Best for: daily coursework, scripting, notebooks, lightweight projects.
\item Strengths: search across files, integrated terminal, extensions, formatting, linting.
\item Trade-offs: can become plugin-heavy; configuration drift across machines.
\item Recommended student default: choose one primary GUI editor and learn it well.
\end{itemize}

\subsection{IDEs (Visual Studio, Eclipse, Xcode, etc.)}
\begin{itemize}
\item Best for: large language ecosystems and complex build systems (Java, C\#, Swift, etc.).
\item Strengths: deep refactoring, project models, integrated debugging/build tooling.
\item Trade-offs: heavier setup; more concepts up front; project configuration can be fragile.
\end{itemize}

\subsection{A decision table (student version)}
\begin{itemize}
\item \textbf{Quick edit on a remote server:} terminal editor.
\item \textbf{Python scripting + small projects:} GUI editor.
\item \textbf{Course uses a specific ecosystem (e.g., Java/C\#/iOS):} IDE.
\item \textbf{When unsure:} GUI editor with an integrated terminal.
\end{itemize}

\section{Essential editor skills (transfer across tools)}
\subsection{Open, save, and "where did it go?"}
\begin{itemize}
\item Always know the \textbf{file path} you are editing.
\item Understand "Save" vs "Save As".
\item Recognize read-only files and permission prompts.
\item Build a habit: confirm the directory and filename after saving.
\end{itemize}

\subsection{Undo/redo and safe experimentation}
\begin{itemize}
\item Undo is your first safety net; version control is your second.
\item Save small, frequent changes; keep commits focused.
\end{itemize}

\subsection{Indentation and whitespace}
\begin{itemize}
\item Choose a consistent indentation style for a project.
\item Configure your editor to show or make visible whitespace when debugging.
\item Understand tabs vs spaces and why Python cares.
\end{itemize}

\subsection{Find, replace, and refactoring safety}
\begin{itemize}
\item Use "Find" before "Replace".
\item Prefer "Replace with confirmation" for nontrivial changes.
\item Learn scope: current selection, current file, entire workspace.
\item Learn regular expressions (basic level): anchors, wildcards, capture groups.
\item After replacing: run tests or re-run the script.
\end{itemize}

\subsection{Multi-file search and navigation}
\begin{itemize}
\item Use workspace search to locate symbols and strings.
\item Learn quick navigation: go to file, go to line, go to definition.
\item Keep a mental model of your project root.
\end{itemize}

\section{Simple scripting workflows in an editor (novice level)}
\subsection{Write-run-read-repeat loop}
\begin{enumerate}
\item Write a small script (10--50 lines).
\item Run it from the integrated terminal.
\item Read the error message carefully.
\item Jump to the referenced file/line.
\item Make one change at a time and re-run.
\end{enumerate}

\subsection{Editor features that help beginners}
\begin{itemize}
\item Syntax highlighting and bracket matching.
\item Auto-indentation and formatting-on-save.
\item Linting messages as "early warnings".
\item Inline documentation (hover tooltips) and jump-to-definition.
\end{itemize}

\subsection{When a full debugger helps}
\begin{itemize}
\item Use a debugger when print/logging is too slow or you need to inspect state.
\item Learn the minimum: breakpoints, step over/into, variable inspection.
\end{itemize}

\section{Debugging files (common student failure modes)}
\subsection{Wrong file, wrong place}
\begin{itemize}
\item Symptom: edits do not change program behavior.
\item Cause: you edited a duplicate file or ran from a different directory.
\item Fix: confirm the working directory; search for duplicates; use absolute paths temporarily.
\end{itemize}

\subsection{Wrong extension or hidden extension}
\begin{itemize}
\item Symptom: file opens in the wrong program or fails to run.
\item Fix: show file extensions; confirm \texttt{.py}, \texttt{.csv}, \texttt{.txt}.
\end{itemize}

\subsection{Encoding and invisible characters}
\begin{itemize}
\item Symptom: weird symbols, parsing errors, or "invalid character" errors.
\item Fix: re-save as UTF-8; inspect for non-printing characters; normalize line endings.
\end{itemize}

\subsection{Line endings (Windows/macOS collaboration)}
\begin{itemize}
\item Symptom: scripts fail, diffs look huge, or tools complain about CRLF/LF.
\item Fix: configure editor to use consistent line endings; rely on repo configuration if provided.
\end{itemize}

\subsection{Config files and indentation-sensitive formats}
\begin{itemize}
\item YAML and similar formats can break with tabs or incorrect indentation.
\item Use visible whitespace; validate with the tool that reads the config.
\end{itemize}

\section{Terminal editors: survival skills}
\subsection{nano essentials (minimum viable)}
\begin{itemize}
\item Open, save, exit without panic.
\item Search and replace.
\item Go to line number.
\end{itemize}

\subsection{vim essentials (minimum viable)}
\begin{itemize}
\item Modes: normal vs insert; how to get back to normal.
\item Save and quit.
\item Search and substitute (with confirmation).
\item Why vim feels fast: motions + operators (optional preview).
\end{itemize}

\subsection{emacs essentials (minimum viable)}
\begin{itemize}
\item Open/save/exit.
\item Search.
\item Basic multi-buffer navigation.
\end{itemize}

\section{GUI editors: practical configuration (keep it minimal)}
\subsection{The minimal settings that matter}
\begin{itemize}
\item Indentation and whitespace display.
\item Format on save (using a formatter appropriate to your language).
\item A linter for early feedback.
\item Integrated terminal.
\end{itemize}

\subsection{Extensions/plugins: a controlled approach}
\begin{itemize}
\item Install only what you can explain.
\item Prefer widely used, well-maintained extensions.
\item Avoid overlapping extensions that fight each other.
\item Document your essential extensions for reproducibility.
\end{itemize}

\subsection{Workspace settings vs global settings}
\begin{itemize}
\item Use workspace settings for project-specific behavior.
\item Keep global settings simple so you can reproduce work on another machine.
\end{itemize}

\section{IDEs: fundamentals without overwhelm}
\subsection{Project model}
\begin{itemize}
\item IDEs often manage build configuration, dependencies, and run targets.
\item Learn where the IDE stores project settings and how to reset or re-import.
\end{itemize}

\subsection{Refactoring and code navigation}
\begin{itemize}
\item IDE refactors are powerful but must be reviewed.
\item Always run tests/build after refactors.
\end{itemize}

\subsection{Debugging integration}
\begin{itemize}
\item IDEs make breakpoints and stepping easy.
\item Learn the minimum and rely on it when print-debugging is not enough.
\end{itemize}

\section{Best practices: habits that prevent pain}
\subsection{Pick a primary editor and a fallback}
\begin{itemize}
\item Primary: your daily GUI editor or required IDE.
\item Fallback: nano (and optionally vim) for remote/emergency edits.
\end{itemize}

\subsection{Treat find/replace as a ``power tool"}
\begin{itemize}
\item Narrow scope.
\item Preview changes.
\item Make a commit before large replacements.
\item Verify with tests or a run.
\end{itemize}

\subsection{Keep text files boring}
\begin{itemize}
\item Use UTF-8.
\item Normalize line endings.
\item Avoid editing generated files by hand.
\item Prefer configuration in version control.
\end{itemize}

\subsection{Use version control as an editor safety net}
\begin{itemize}
\item Commit logical steps.
\item Use diffs to review changes before running.
\item Revert when experiments go wrong.
\end{itemize}

\section{Worked examples (outline)}
\subsection{Example 1: Write and run a tiny script}
\begin{itemize}
\item Create \texttt{hello.py}.
\item Run from the terminal.
\item Fix a syntax error using the editor's line/column cues.
\end{itemize}

\subsection{Example 2: Debug a broken config file}
\begin{itemize}
\item Identify indentation and invisible-character issues.
\item Turn on visible whitespace.
\item Validate by re-running the tool.
\end{itemize}

\subsection{Example 3: Safe rename/refactor with find/replace}
\begin{itemize}
\item Use multi-file search.
\item Replace with confirmation.
\item Run tests.
\end{itemize}

\subsection{Example 4: Remote emergency edit over SSH}
\begin{itemize}
\item Open a config file in nano.
\item Make a minimal change.
\item Save/exit; verify.
\end{itemize}

\section{Templates}
\subsection{Template A: Editor setup checklist (first week)}
\begin{verbatim}
Editor:

* Installed and updated
* Default indentation configured
* Visible whitespace toggle known
* Format-on-save configured (if course uses it)
* Linter installed (if course uses it)
* Integrated terminal enabled
  Fallback editor:
* nano basics practiced (open/save/exit/search)
  \end{verbatim}

\subsection{Template B: Safe find/replace protocol}
\begin{verbatim}

1. Search (no replace) and inspect matches
2. Narrow scope (file/selection/project)
3. Replace with confirmation
4. Save and review diff
5. Run tests or rerun script
6. Commit with a message describing the change
   \end{verbatim}

\section{Exercises}
\begin{enumerate}
\item Configure your editor to show line numbers and visible whitespace; explain what you changed.
\item Write a short script and fix two intentional errors using editor cues.
\item Perform a multi-file search in a small project; report where a function name appears.
\item Run a safe find/replace with confirmation and verify by running tests.
\item Open and edit a file in nano; save and exit confidently.
\item Optional: learn a minimal vim workflow (insert, save, quit, search).
\end{enumerate}

\section{One-page checklist}
\begin{itemize}
\item I can open, save, and locate files reliably.
\item I can use find/replace safely, including across multiple files.
\item I can debug common file issues (paths, extensions, encoding, line endings).
\item I understand when to use terminal editors vs GUI editors vs IDEs.
\item My editor is minimally configured (indentation, formatting, linting).
\item I use version control to review and recover from editing mistakes.
\item I have a fallback editor skill for remote/emergency use.
\end{itemize}

\appendix
\section{Quick reference: terminal editor survival commands (optional handout)}
\subsection{nano}
\begin{itemize}
\item Open: \texttt{nano file}
\item Save: (document the keystroke)
\item Exit: (document the keystroke)
\item Find/replace: (document the keystroke)
\end{itemize}

\subsection{vim}
\begin{itemize}
\item Insert mode: \texttt{i}
\item Normal mode: \texttt{Esc}
\item Save/quit: \texttt{:wq}
\item Quit without saving: \texttt{:q!}
\item Search: \texttt{/pattern}
\item Substitute (conceptual): \texttt{:%s/old/new/gc}
\end{itemize}

\section{Quick reference: GUI/IDE search}
\begin{itemize}
\item Find in file, replace in file
\item Find in workspace, replace in workspace
\item Go to line, go to file, go to definition
\end{itemize}
